% GENERATED FILE. Edit canonical lessons, then run npm run generate:books.
% canonical-source-hash: fnv1a64:8ca580afa220be23
% canonical-lessons: RU-C418-peresazhivatsya, RU-C418-konechnaya, RU-C418-prigorodnyy, RU-C418-proyezzhat, RU-C418-provozhat, RU-R418-review-411-418

\chapter{Changing trains}
\label{ch:ru-changing-trains}
\hlchaptermodality{418}

\begin{chapteropening}
\textbf{By the end of this chapter:} \emph{I can travel with a change: changing trains, the last stop, local trains, going past and seeing someone off.}

Complete the last lesson of chapter 418: I can travel with a change: changing trains, the last stop, local trains, going past and seeing someone off.
\end{chapteropening}

\section[\texorpdfstring{\ru{пересаживаться} (peresázhivatsya)}{пересаживаться (peresázhivatsya)}]{\ru{пересаживаться} (peresázhivatsya) — to change trains or buses}

\label{lesson:RU-C418-peresazhivatsya}



\begin{quote}
Before the new one: say the Russian for a route, then the Russian for a fixed-route minibus.
\end{quote}

\subsection*{You'll want to know: \ru{пересаживаться}}
\textbf{\ru{пересаживаться}} (\emph{peresázhivatsya}) — \textquotedblleft{}to change trains or buses\textquotedblright{}.

\begin{quote}
\ru{Мы} \ru{пересаживаемся} \ru{на} \ru{другой} \ru{поезд} \ru{в} \ru{центре}.
\end{quote}

\begin{quote}
\emph{My peresázhivayemsya na drugóy póyezd v tséntre.}
\end{quote}

\begin{quote}
We change to another train in the centre.
\end{quote}

Imperfective; perfective \emph{\ru{пересесть}} (\emph{\ru{пересяду}, \ru{пересядешь}}). Takes \emph{\ru{на}} + accusative.

\begin{cousinweb}
\emph{\ru{пересаживаться}} is built from \emph{\ru{пере}-} 'across', \emph{-\ru{саживать}}, from \emph{\ru{садить}} 'to seat', and \emph{-\ru{ся}}: to seat oneself somewhere else. EtymDB lists \emph{\ru{садить}} without a deeper source.
\end{cousinweb}

\begin{grammarlens}[title={What you've built}]
One more word for telling, explaining and describing at B1.
\end{grammarlens}

\subsection*{Your turn}
\begin{itemize}
\raggedright
  \item \emph{Say it:} \emph{peresázhivatsya}
  \item \emph{Say it:} \emph{peresázhivatsya}, once more
  \item \emph{Say it:} a sentence of your own with \emph{peresázhivatsya}
  \item \emph{Recall:} say the Russian for a route, then the Russian for a fixed-route minibus, then say \emph{peresázhivatsya} again
\end{itemize}

\begin{tcolorbox}[breakable,colback=teal!4,colframe=teal!35!black,title={Before you move on}]
What does \ru{пересаживаться} mean? (\textquotedblleft{}To change trains or buses\textquotedblright{}.) Say it once more.
\end{tcolorbox}
\section[\texorpdfstring{\ru{конечная} (konéchnaya)}{конечная (konéchnaya)}]{\ru{конечная} (konéchnaya) — the last stop, the terminus}

\label{lesson:RU-C418-konechnaya}



\begin{quote}
Before the new one: say the Russian for a fixed-route minibus, then the Russian for to change trains or buses.
\end{quote}

\subsection*{You'll want to know: \ru{конечная}}
\textbf{\ru{конечная}} (\emph{konéchnaya}) — \textquotedblleft{}the last stop, the terminus\textquotedblright{}.

\begin{quote}
\ru{Мы} \ru{едем} \ru{до} \ru{конечной}.
\end{quote}

\begin{quote}
\emph{My yédem do konéchnoy.}
\end{quote}

\begin{quote}
We are going to the last stop.
\end{quote}

Short for \emph{\ru{конечная} \ru{остановка}}; an adjective used as a noun.

\begin{cousinweb}
\emph{\ru{конечная}} is short for \emph{\ru{конечная} \ru{остановка}} 'the final stop'. \emph{\ru{Конечный}} is built on \emph{\ru{конец}} 'an end', which is inherited from Proto-Slavic \emph{konĭcĭ} 'an end' (EtymDB).
\end{cousinweb}

\begin{grammarlens}[title={What you've built}]
One more word for telling, explaining and describing at B1.
\end{grammarlens}

\subsection*{Your turn}
\begin{itemize}
\raggedright
  \item \emph{Say it:} \emph{konéchnaya}
  \item \emph{Say it:} \emph{konéchnaya}, once more
  \item \emph{Say it:} a sentence of your own with \emph{konéchnaya}
  \item \emph{Recall:} say the Russian for a fixed-route minibus, then the Russian for to change trains or buses, then say \emph{konéchnaya} again
\end{itemize}

\begin{tcolorbox}[breakable,colback=teal!4,colframe=teal!35!black,title={Before you move on}]
What does \ru{конечная} mean? (\textquotedblleft{}The last stop, the terminus\textquotedblright{}.) Say it once more.
\end{tcolorbox}
\section[\texorpdfstring{\ru{пригородный} (prígorodnyy)}{пригородный (prígorodnyy)}]{\ru{пригородный} (prígorodnyy) — suburban, local (of trains and buses)}

\label{lesson:RU-C418-prigorodnyy}



\begin{quote}
Before the new one: say the Russian for to change trains or buses, then the Russian for the last stop.
\end{quote}

\subsection*{You'll want to know: \ru{пригородный}}
\textbf{\ru{пригородный}} (\emph{prígorodnyy}) — \textquotedblleft{}suburban, local (of trains and buses)\textquotedblright{}.

\begin{quote}
\ru{Бабушка} \ru{едет} \ru{на} \ru{дачу} \ru{на} \ru{пригородном} \ru{поезде}.
\end{quote}

\begin{quote}
\emph{Bábushka yédet na dáchu na prígorodnom póyezde.}
\end{quote}

\begin{quote}
Grandma goes to her country cottage on the local train.
\end{quote}

The stress is on the first syllable: \emph{prígorodnyy}.

\begin{cousinweb}
\emph{\ru{пригородный}} is built from \emph{\ru{при}-} 'near', \emph{\ru{город}} 'a town' and \emph{-\ru{ный}}. \emph{\ru{Город}} is inherited from Proto-Slavic \emph{gordŭ} 'an enclosure' (EtymDB).
\end{cousinweb}

\begin{grammarlens}[title={What you've built}]
One more word for telling, explaining and describing at B1.
\end{grammarlens}

\subsection*{Your turn}
\begin{itemize}
\raggedright
  \item \emph{Say it:} \emph{prígorodnyy}
  \item \emph{Say it:} \emph{prígorodnyy}, once more
  \item \emph{Say it:} a sentence of your own with \emph{prígorodnyy}
  \item \emph{Recall:} say the Russian for to change trains or buses, then the Russian for the last stop, then say \emph{prígorodnyy} again
\end{itemize}

\begin{tcolorbox}[breakable,colback=teal!4,colframe=teal!35!black,title={Before you move on}]
What does \ru{пригородный} mean? (\textquotedblleft{}Suburban, local (of trains and buses)\textquotedblright{}.) Say it once more.
\end{tcolorbox}
\section[\texorpdfstring{\ru{проезжать} (proyezzhát)}{проезжать (proyezzhát)}]{\ru{проезжать} (proyezzhát) — to go past, to drive through; to miss (one's stop)}

\label{lesson:RU-C418-proyezzhat}



\begin{quote}
Before the new one: say the Russian for the last stop, then the Russian for suburban.
\end{quote}

\subsection*{You'll want to know: \ru{проезжать}}
\textbf{\ru{проезжать}} (\emph{proyezzhát}) — \textquotedblleft{}to go past, to drive through; to miss (one's stop)\textquotedblright{}.

\begin{quote}
\ru{Автобус} \ru{проезжает} \ru{мимо} \ru{нашего} \ru{дома}.
\end{quote}

\begin{quote}
\emph{Avtóbus proyezzháyet mímo náshego dóma.}
\end{quote}

\begin{quote}
The bus goes past our house.
\end{quote}

Imperfective; perfective \emph{\ru{проехать}} (\emph{\ru{проеду}, \ru{проедешь}}). \emph{\ru{Я} \ru{проехал} \ru{свою} \ru{остановку}} — I missed my stop.

\begin{cousinweb}
\emph{\ru{проезжать}} is built from \emph{\ru{про}-} 'past, through' and \emph{-\ru{езжать}}, from the root of \emph{\ru{ездить}} 'to ride, to travel'. EtymDB lists \emph{\ru{проезжать}} without a deeper source.
\end{cousinweb}

\begin{grammarlens}[title={What you've built}]
One more word for telling, explaining and describing at B1.
\end{grammarlens}

\subsection*{Your turn}
\begin{itemize}
\raggedright
  \item \emph{Say it:} \emph{proyezzhát}
  \item \emph{Say it:} \emph{proyezzhát}, once more
  \item \emph{Say it:} a sentence of your own with \emph{proyezzhát}
  \item \emph{Recall:} say the Russian for the last stop, then the Russian for suburban, then say \emph{proyezzhát} again
\end{itemize}

\begin{tcolorbox}[breakable,colback=teal!4,colframe=teal!35!black,title={Before you move on}]
What does \ru{проезжать} mean? (\textquotedblleft{}To go past, to drive through; to miss (one's stop)\textquotedblright{}.) Say it once more.
\end{tcolorbox}
\section[\texorpdfstring{\ru{провожать} (provozhát)}{провожать (provozhát)}]{\ru{провожать} (provozhát) — to see off}

\label{lesson:RU-C418-provozhat}



\begin{quote}
Before the new one: say the Russian for suburban, then the Russian for to go past.
\end{quote}

\subsection*{You'll want to know: \ru{провожать}}
\textbf{\ru{провожать}} (\emph{provozhát}) — \textquotedblleft{}to see off\textquotedblright{}.

\begin{quote}
\ru{Родители} \ru{провожают} \ru{сына} \ru{на} \ru{вокзале}.
\end{quote}

\begin{quote}
\emph{Rodíteli provozháyut sýna na vokzále.}
\end{quote}

\begin{quote}
The parents are seeing their son off at the station.
\end{quote}

Imperfective; perfective \emph{\ru{проводить}} (\emph{\ru{провожу}, \ru{проводишь}}). The imperfective \emph{\ru{проводить}} is a different verb — to spend (time).

\begin{cousinweb}
\emph{\ru{провожать}} is built on \emph{\ru{проводить}} 'to lead along, to accompany', from \emph{\ru{про}-} and \emph{\ru{водить}} 'to lead'. \emph{\ru{Водить}} is inherited from Proto-Slavic \emph{voditi} (EtymDB).
\end{cousinweb}

\begin{grammarlens}[title={What you've built}]
One more word for telling, explaining and describing at B1.
\end{grammarlens}

\subsection*{Your turn}
\begin{itemize}
\raggedright
  \item \emph{Say it:} \emph{provozhát}
  \item \emph{Say it:} \emph{provozhát}, once more
  \item \emph{Say it:} a sentence of your own with \emph{provozhát}
  \item \emph{Recall:} say the Russian for suburban, then the Russian for to go past, then say \emph{provozhát} again
\end{itemize}

\begin{tcolorbox}[breakable,colback=teal!4,colframe=teal!35!black,title={Before you move on}]
What does \ru{провожать} mean? (\textquotedblleft{}To see off\textquotedblright{}.) Say it once more.
\end{tcolorbox}
\section[(dialogue)]{Words from chapters 411–418 — a review}

\label{lesson:RU-R418-review-411-418}



\begin{quote}
Say the Russian for a fellow traveller and a lift.
\end{quote}

\subsection*{Your turn}
Say these aloud:

\begin{itemize}
\raggedright
  \item a fellow traveller, a lift, to give, to drop off, to fasten one's seat belt
  \item a pedestrian, a cyclist, to give way, to let through, road markings
  \item a fare, an open sleeping car, a compartment, a rate, a concession
  \item a station platform, a platform, a departures board, departure, arrival
  \item to set off, to arrive, to pull out, to announce, to make it
  \item a sleeping-car attendant, a vestibule, a porter, a trolley, to present
  \item a turnstile, a travel pass, a fare dodger, a route, a fixed-route minibus
  \item to change trains or buses, the last stop, suburban, to go past, to see off
\end{itemize}

\begin{tcolorbox}[breakable,colback=teal!4,colframe=teal!35!black,title={Before you move on}]
A fellow traveller? (\textbf{\ru{попутчик}}.) A lift? (\textbf{\ru{попутка}}.) To give? (\textbf{\ru{подвозить}}.) To drop off? (\textbf{\ru{высаживать}}.) To fasten one's seat belt? (\textbf{\ru{пристёгиваться}}.) A pedestrian? (\textbf{\ru{пешеход}}.) A cyclist? (\textbf{\ru{велосипедист}}.) To give way? (\textbf{\ru{уступать}}.) To let through? (\textbf{\ru{пропускать}}.) Road markings? (\textbf{\ru{разметка}}.) A fare? (\textbf{\ru{проезд}}.) An open sleeping car? (\textbf{\ru{плацкарт}}.) A compartment? (\textbf{\ru{купе}}.) A rate? (\textbf{\ru{тариф}}.) A concession? (\textbf{\ru{льгота}}.) A station platform? (\textbf{\ru{перрон}}.) A platform? (\textbf{\ru{платформа}}.) A departures board? (\textbf{\ru{табло}}.) Departure? (\textbf{\ru{отправление}}.) Arrival? (\textbf{\ru{прибытие}}.) To set off? (\textbf{\ru{отправляться}}.) To arrive? (\textbf{\ru{прибывать}}.) To pull out? (\textbf{\ru{отходить}}.) To announce? (\textbf{\ru{объявлять}}.) To make it? (\textbf{\ru{успевать}}.) A sleeping-car attendant? (\textbf{\ru{проводник}}.) A vestibule? (\textbf{\ru{тамбур}}.) A porter? (\textbf{\ru{носильщик}}.) A trolley? (\textbf{\ru{тележка}}.) To present? (\textbf{\ru{предъявлять}}.) A turnstile? (\textbf{\ru{турникет}}.) A travel pass? (\textbf{\ru{проездной}}.) A fare dodger? (\textbf{\ru{безбилетник}}.) A route? (\textbf{\ru{маршрут}}.) A fixed-route minibus? (\textbf{\ru{маршрутка}}.) To change trains or buses? (\textbf{\ru{пересаживаться}}.) The last stop? (\textbf{\ru{конечная}}.) Suburban? (\textbf{\ru{пригородный}}.) To go past? (\textbf{\ru{проезжать}}.) To see off? (\textbf{\ru{провожать}}.)
\end{tcolorbox}
